\chapter{Marchés de l'électricité II~: négociation}\label{ch:m3:power-markets-trading}

À quarante-cinq minutes de seize heures, la prévision d'un parc éolien pour l'heure suivante baisse d'un tiers.
Il avait vendu 100 mégawattheures pour cette heure lors de l'enchère de la veille~; il prévoit désormais d'en produire 70.
Il peut acheter les 30 manquants sur le \omterm{def:m3:power-markets-trading:intraday}{marché infrajournalier} continu, au prix qu'affiche le carnet d'ordres,
ou livrer en déficit et laisser le gestionnaire du réseau lui facturer le \omterm{def:m3:power-markets-trading:imbalance}{prix de règlement des écarts},
qu'il ne connaîtra que quelques semaines plus tard. La plupart du temps, l'écart ne représente que quelques euros.
Lors d'une soirée tendue, il peut atteindre des centaines d'euros. La négociation d'électricité après l'enchère du \omterm{def:m3:power-markets-design:da}{marché J-1}
consiste à gérer cet écart~: des prévisions qui évoluent, un marché qui ferme peu avant la livraison,
et un règlement qui sanctionne celui qui se trompe. Ce chapitre traite du \omterm{def:m3:power-markets-trading:intraday}{marché infrajournalier}, du règlement des écarts,
des instruments de couverture de la congestion, de l'économie de la production renouvelable
et de ses contrats de long terme, ainsi que des batteries.

\section{Marchés infrajournaliers continus}

\begin{definition}[Marché infrajournalier, heure de fermeture du guichet]\label{def:m3:power-markets-trading:intraday}
Un \emph{marché infrajournalier}\index{marché infrajournalier} négocie de l'électricité pour une livraison plus tard le jour même,
ou le lendemain, après l'enchère du \omterm{def:m3:power-markets-design:da}{marché J-1}, soit en continu dans un carnet d'ordres
à priorité prix-temps (livre~1, chapitre~19), soit lors d'enchères infrajournalières. Son \emph{heure de
fermeture du guichet}\index{heure de fermeture du guichet} pour une MTU est le dernier instant auquel cette MTU peut être négociée~;
au-delà, les écarts sont réglés comme des déséquilibres.
\end{definition}

Le négoce infrajournalier existe parce que les prévisions s'améliorent à mesure que la livraison approche~: une prévision éolienne établie
un jour à l'avance peut s'écarter d'un cinquième de la production, contre beaucoup moins une heure à l'avance. Chaque producteur et chaque fournisseur
corrige sa position au fil de l'évolution de sa prévision, et les prix évoluent avec l'erreur de prévision agrégée
de la zone.

\begin{dated}{2026-09}{Couplage infrajournalier européen}\label{dat:m3:power-markets-trading:sidc}
Le Single Intraday Coupling (SIDC), lancé les 12 et 13 juin 2018 dans quinze pays,
permet aux ordres d'une zone de se confronter à ceux d'une autre tant que la capacité transfrontalière est disponible. Les échanges
transfrontaliers se fermaient 60 minutes avant le début de chaque MTU~; une fermeture du guichet transfrontalière à 30 minutes
est depuis entrée en service sur certaines frontières. Les échanges au sein d'une zone peuvent se poursuivre plus près de la livraison,
selon le pays.
\end{dated}

\begin{omfigure}
\centering
\begin{tikzpicture}[font=\small]
\draw[thick,-{Stealth[length=2.5mm]}] (0,0) -- (13.2,0);
\foreach \x/\t/\l in {0.4/{des mois\\avant}/{forwards et\\contrats à terme},3.2/{J$-1$, 12:00}/{enchère\\du jour J-1},6.4/{de J$-1$ après-midi\\à la livraison}/{infrajournalier\\continu},9.4/{fermeture\\du guichet}/{dernière\\transaction},12.0/{des semaines\\après}/{règlement\\des déséquilibres}}{
  \draw[thick] (\x,0.12) -- (\x,-0.12);
  \node[anchor=south,align=center] at (\x,0.2) {\t};
  \node[anchor=north,align=center] at (\x,-0.2) {\l};}
\draw[omThm,thick] (10.7,0.3) -- (10.7,-0.3);
\node[omThm,anchor=north,font=\footnotesize] at (10.7,-1.45) {livraison};
\end{tikzpicture}

\omcaption{Les marchés d'une MTU, depuis le marché à terme jusqu'au règlement. Chaque transaction conclue avant
la fermeture du guichet modifie le programme~; chaque écart constaté après celle-ci constitue un déséquilibre. Schéma.}\label{fig:m3:power-markets-trading:timeline}
\end{omfigure}

\section{Règlement des écarts et responsabilité d'équilibre}

\begin{definition}[Responsable d'équilibre]\label{def:m3:power-markets-trading:brp}
Un \emph{responsable d'équilibre}\index{responsable d'équilibre} (balancing responsible party, BRP) est un intervenant de marché,
ou son représentant, financièrement responsable de l'écart entre
l'énergie qu'il a programmée et l'énergie qu'il a effectivement injectée ou soutirée, pour un portefeuille de
points de raccordement, à chaque période de règlement des écarts.
\end{definition}

\begin{definition}[Volume d'écart, prix de règlement des écarts]\label{def:m3:power-markets-trading:imbalance}
Le \emph{volume d'écart}\index{volume d'écart} d'un \omterm{def:m3:power-markets-trading:brp}{responsable d'équilibre} pour une période de règlement est son
injection nette mesurée moins son injection nette programmée~: positif lorsqu'il est en excédent (il a livré
davantage qu'il n'a vendu), négatif lorsqu'il est en déficit. Le \emph{prix de règlement des écarts}\index{prix de règlement des écarts} est
le prix auquel le gestionnaire du réseau le règle, déterminé à partir du coût de l'énergie d'équilibrage
que ce gestionnaire a activée au cours de cette période.
\end{definition}

Dans un système à prix unique, un \omterm{def:m3:power-markets-trading:brp}{responsable d'équilibre} en excédent reçoit le \omterm{def:m3:power-markets-trading:imbalance}{prix de règlement des écarts} pour son surplus, et un responsable en déficit
le paie pour son manque. Lorsque l'ensemble du système est en déficit, le gestionnaire active des réserves à la hausse
et le \omterm{def:m3:power-markets-trading:imbalance}{prix de règlement des écarts} s'élève~: être en déficit coûte alors cher, et être en excédent est
récompensé. Le \omterm{def:m3:power-markets-trading:imbalance}{prix de règlement des écarts} n'est donc une pénalité que pour ceux qui vont dans le même
sens que le système, et un \omterm{def:m3:power-markets-trading:brp}{responsable d'équilibre} dont l'erreur est de signe opposé en tire un gain. La plupart
des systèmes européens règlent les écarts par tranches de 15 minutes~; l'Allemagne calcule un prix uniforme par
quart d'heure pour l'ensemble de ses groupes d'équilibrage.

\begin{example}[L'après-midi d'un parc éolien]\label{ex:m3:power-markets-trading:wind}
Un parc éolien vend 80, 90, 100 et 100~MWh en J-1 pour quatre heures, à 62, 58, 55 et
\qty{57}{EUR/MWh}. Avant la dernière heure, sa prévision baisse de 30~MWh et il rachète 30 en infrajournalier à
\qty{170}{EUR/MWh}. Il mesure 82, 88, 101 et 55~MWh~; les \omterm{def:m3:power-markets-trading:imbalance}{prix de règlement des écarts} sont de 60,
75, 50 et \qty{240}{EUR/MWh}. Son revenu de négoce s'élève à EUR~16\,280 (les ventes en J-1 moins
le rachat), et ses écarts ($+2$, $-2$, $+1$ et $-15$~MWh) lui coûtent EUR~3\,580, presque
tous lors de la dernière heure, où il était en déficit alors que le système l'était aussi~: EUR~12\,700 au total.
\end{example}

\section{Droits de transport et offres virtuelles}

Dans les marchés nodaux, un producteur est payé au prix de son nœud et un fournisseur paie au prix de son nœud de consommation, si bien que
la congestion entre les deux constitue un risque que le marché leur permet de couvrir.

\begin{definition}[Droit financier de transport]\label{def:m3:power-markets-trading:ftr}
Un \emph{droit financier de transport}\index{droit financier de transport} (FTR) du nœud $A$
vers le nœud $B$ donne à son détenteur, pour chaque heure de sa durée, un droit égal à une quantité multipliée par le prix de congestion du \omterm{def:m3:power-markets-design:da}{marché J-1}
au nœud $B$ moins celui du nœud $A$~: un paiement lorsque le chemin est congestionné dans le
sens du droit et, pour une obligation, une charge lorsqu'il est congestionné dans l'autre sens.
Il est vendu par le gestionnaire du réseau lors d'enchères et financé par les \omterm{def:m3:power-markets-design:lmp}{rentes de congestion} qu'il perçoit.
\end{definition}

Un producteur situé en $A$ qui vend à un consommateur situé en $B$ à un prix fixe supporte le risque de congestion entre
les deux~; un FTR du nœud $A$ vers le nœud $B$ le transforme en un prix d'enchère fixe.

\begin{definition}[Offre virtuelle]\label{def:m3:power-markets-trading:virtual}
Une \emph{offre virtuelle}\index{offre virtuelle} est une position purement financière dans un marché à double règlement
~: une offre de vente (offre à la hausse) ou une offre d'achat (offre à la baisse) à un nœud sur le
\omterm{def:m3:power-markets-design:da}{marché J-1}, liquidée automatiquement au prix temps réel, sans livraison
physique.
\end{definition}

Un trader qui anticipe des prix temps réel supérieurs aux prix J-1 à un nœud achète en J-1 au moyen
d'une offre d'achat à la baisse et perçoit la différence~; le marché y trouve son intérêt, car cela rapproche le
prix J-1 du prix temps réel attendu. Les mêmes transactions peuvent servir à manipuler~: une
\omterm{def:m3:power-markets-trading:virtual}{offre virtuelle} qui perd de l'argent par elle-même mais fait monter la valeur des FTR du trader sur le même chemin
constitue l'abus classique, et les règles de PJM suspendent alors le versement des paiements de FTR dans de tels cas.

\section{Production renouvelable, prix captés et contrats d'achat d'électricité}

Un parc éolien ou un parc solaire ne choisit pas le moment où il produit. Ce qu'il gagne sur le marché spot
n'est pas le prix moyen, mais la moyenne pondérée par sa propre production.

\begin{definition}[Prix capté]\label{def:m3:power-markets-trading:capture}
Le \emph{prix capté}\index{prix capté} d'un producteur sur une période est la moyenne
du prix de marché pondéré par sa production, $\sum_h p_h q_h / \sum_h q_h$~; son taux de capture est
le prix capté divisé par le prix moyen pondéré dans le temps (prix de base).
\end{definition}

Les parcs solaires d'une même zone produisent tous aux mêmes heures, celles où ils font
baisser le prix (\cref{ch:m3:power-markets-design})~: plus on construit de solaire, plus son taux de capture
baisse. Dans les données allemandes, le \omterm{def:m3:power-markets-trading:capture}{prix capté} par le solaire s'élève à \qty{46.23}{EUR/MWh} en 2024, contre une
moyenne de base de \qty{78.51}{} (un taux de capture de 59~\%), et \qty{46.07}{} contre
\qty{89.32}{EUR/MWh} en 2025 (52~\%)~; l'éolien terrestre capte 82~\% et 87~\%
(\cref{fig:m3:power-markets-trading:capture}).

\begin{omfigure}
\begin{tikzpicture}
\begin{axis}[width=11cm,height=4.8cm,
  ylabel={EUR/MWh},
  tick label style={font=\small},label style={font=\small},
  xmin=2024,xmax=2026,ymin=0,ymax=160,grid=major,grid style={gray!25},
  xtick={2024,2024.5,2025,2025.5,2026},xticklabels={janv. 2024,juil. 2024,janv. 2025,juil. 2025,janv. 2026},
  legend columns=2,legend style={font=\footnotesize,at={(0.5,-0.2)},anchor=north,draw=none,fill=none,/tikz/every even column/.append style={column sep=8pt}},
  table/col sep=comma]
\addplot[omDef,very thick,mark=*,mark size=1.2pt] table[x=t,y=base]{figdata/markets-3/06-power-markets-trading/capture.csv};
\addplot[omProp,very thick,mark=square*,mark size=1.2pt] table[x=t,y=solar]{figdata/markets-3/06-power-markets-trading/capture.csv};
\legend{moyenne de base,prix capté du solaire}
\end{axis}
\end{tikzpicture}

\omcaption{Allemagne, par mois~: le prix moyen du \omterm{def:m3:power-markets-design:da}{marché J-1} et le \omterm{def:m3:power-markets-trading:capture}{prix capté} par la production
solaire, 2024--2025. L'écart est le plus large dans les mois les plus ensoleillés. Données~: Bundesnetzagentur |
SMARD.de, CC BY 4.0.}\label{fig:m3:power-markets-trading:capture}
\end{omfigure}

\begin{definition}[Contrat d'achat d'électricité, risque de profil]\label{def:m3:power-markets-trading:ppa}
Un \emph{contrat d'achat d'électricité}\index{contrat d'achat d'électricité} (PPA) est un contrat de long terme
par lequel un acheteur, souvent un grand consommateur, achète la production d'un producteur (ou un équivalent
financier) à un prix convenu pour plusieurs années. Le \emph{risque de profil}\index{risque de profil} est le risque
que le profil de la production livrée a une valeur différente de celle d'un bloc de base de même
volume~: dans un PPA à production variable, il pèse sur l'acheteur~; dans un PPA de base, sur le vendeur, qui
doit acheter aux heures où il ne produit pas.
\end{definition}

Un PPA solaire à production variable conclu à \qty{60}{EUR/MWh} n'était une bonne couverture pour un acheteur en 2024 que s'il
valorisait la production à son \omterm{def:m3:power-markets-trading:capture}{prix capté}~: à un \omterm{def:m3:power-markets-trading:capture}{prix capté} de 46, l'acheteur a payé environ
14~EUR par MWh au-dessus de la valeur de marché de ce qu'il a reçu.

\begin{dated}{2026-09}{Volumes contractés et stockés}\label{dat:m3:power-markets-trading:ppa}
Pexapark a recensé environ 6.08~GW de capacités renouvelables contractées sous forme de PPA en Europe au premier semestre 2025, dans 124 contrats, soit 26~\% de moins
qu'un an plus tôt. En Texas, une présentation devant le Technology and Security Committee de l'ERCOT en février 2026 évaluait la capacité installée et en service des batteries à
15\,712~MW, dont plus de 6\,000~MW ajoutés au cours de l'année précédente.
\end{dated}

\section{Batteries et actifs flexibles}

\begin{definition}[Rendement aller-retour]\label{def:m3:power-markets-trading:rte}
Le \emph{rendement aller-retour}\index{rendement aller-retour} d'un actif de stockage est l'énergie
restituée divisée par l'énergie consommée pour la charger~; le reste est perdu lors de la conversion.
\end{definition}

Une batterie achète les heures bon marché et vend les heures chères. Avec une anticipation parfaite des prix du \omterm{def:m3:power-markets-design:da}{marché J-1},
sa meilleure planification quotidienne est un petit programme dynamique sur son état de charge, un problème du
genre de ceux que traite en général le livre~4, chapitre~9.

\begin{method}[Arbitrage J-1 d'une batterie]\label{met:m3:power-markets-trading:battery}
Discrétiser l'état de charge par pas d'une heure à pleine puissance. Pour chaque heure, à partir de la
première, conserver pour chaque état accessible (niveau de charge, nombre de cycles utilisés) le meilleur gain cumulé jusqu'ici~; on passe ensuite à l'état suivant
en chargeant (payer le prix pour $1/\sqrt{\eta}$ fois l'énergie stockée), en déchargeant (recevoir le
prix pour $\sqrt{\eta}$ fois l'énergie restituée) ou en restant inactif~; à la fin de la journée, retenir le
meilleur état vide.
\end{method}

Une batterie de 100~MW et 200~MWh, avec un \omterm{def:m3:power-markets-trading:rte}{rendement aller-retour} de 88~\% et un cycle complet par jour
(chiffres illustratifs) aurait gagné, avec une anticipation parfaite des prix allemands du
\omterm{def:m3:power-markets-design:da}{marché J-1}, 65\,968~EUR par MW en 2024 et 74\,546~EUR par MW en 2025~; avec deux cycles par jour,
le gain 2025 atteint 88\,688~EUR. Une anticipation parfaite surestime ce qu'un opérateur réel gagne sur la
seule enchère, alors que les \omterm{def:m3:power-markets-trading:intraday}{marchés infrajournaliers} et les réserves ajoutent des revenus que ce calcul omet.

\begin{omfigure}
\begin{tikzpicture}
\begin{axis}[width=11cm,height=4.6cm,axis y line*=right,axis x line=none,
  ylabel={état de charge, MWh},ylabel near ticks,
  tick label style={font=\small},label style={font=\small},
  xmin=-0.5,xmax=23.5,ymin=0,ymax=440,ytick={0,100,200},
  table/col sep=comma]
\addplot[omDef,thick,const plot,fill=omDef!12] table[x=tend,y=soc]{figdata/markets-3/06-power-markets-trading/batteryday.csv}\closedcycle;\label{plot:m3:battsoc}
\end{axis}
\begin{axis}[width=11cm,height=4.6cm,axis y line*=left,
  xlabel={heure, heure locale allemande},ylabel={EUR/MWh},
  tick label style={font=\small},label style={font=\small},
  xmin=-0.5,xmax=23.5,ymin=-20,ymax=200,grid=major,grid style={gray!25},
  xtick={0,3,6,9,12,15,18,21},
  legend columns=2,legend style={font=\footnotesize,at={(0.5,-0.33)},anchor=north,draw=none,fill=none,/tikz/every even column/.append style={column sep=8pt}},
  table/col sep=comma]
\addplot[omThm,very thick,const plot mark mid] table[x=hour,y=price]{figdata/markets-3/06-power-markets-trading/batteryday.csv};\addlegendentry{prix du marché J-1}
\addlegendimage{/pgfplots/refstyle=plot:m3:battsoc}\addlegendentry{état de charge (fin de l'heure)}
\end{axis}
\end{tikzpicture}

\omcaption{La meilleure planification de la batterie sur le \omterm{def:m3:power-markets-design:da}{marché J-1} le 18 juin 2025~: elle charge pendant les deux
heures les moins chères après-midi (prix légèrement négatifs) et décharge pendant les deux plus chères du
soir, gagnant 30\,149~EUR. Batterie illustrative~; prix~: Bundesnetzagentur | SMARD.de,
CC BY 4.0.}\label{fig:m3:power-markets-trading:battery}
\end{omfigure}

\begin{omfigure}
\begin{tikzpicture}
\begin{axis}[width=11cm,height=4.6cm,
  xlabel={jour de l'année},ylabel={milliers d'euros par MW},
  tick label style={font=\small},label style={font=\small},
  xmin=1,xmax=366,ymin=0,ymax=80,grid=major,grid style={gray!25},
  legend columns=2,legend style={font=\footnotesize,at={(0.5,-0.3)},anchor=north,draw=none,fill=none,/tikz/every even column/.append style={column sep=8pt}},
  table/col sep=comma]
\addplot[omProp,very thick] table[x=doy,y=y2024]{figdata/markets-3/06-power-markets-trading/batterycum.csv};
\addplot[omDef,very thick] table[x=doy,y=y2025]{figdata/markets-3/06-power-markets-trading/batterycum.csv};
\legend{2024,2025}
\end{axis}
\end{tikzpicture}

\omcaption{Revenu cumulé de l'arbitrage sur le \omterm{def:m3:power-markets-design:da}{marché J-1} de la batterie illustrative (un cycle par jour,
prévision parfaite), par MW de puissance, à partir du 1\up{er} janvier. Le revenu s'accumule le plus vite
pendant les mois ensoleillés et volatils. Prix~: Bundesnetzagentur | SMARD.de, CC BY 4.0.}\label{fig:m3:power-markets-trading:cum}
\end{omfigure}

\section{Tutoriel~: règlement d'une journée et conduite d'une batterie}\label{tut:m3:power-markets-trading:battery}

\begin{tutorial}
\textbf{Objectif.} Régler la journée d'un producteur et programmer une batterie face aux prix J-1.
\textbf{État final~:} \cref{ex:m3:power-markets-trading:wind},
\cref{fig:m3:power-markets-trading:battery,fig:m3:power-markets-trading:cum} et les revenus
annuels par MW.
\begin{tutsteps}
\item \textbf{Règlement et \omterm{def:m3:power-markets-trading:capture}{prix capté}.}
\omcode{code/firm/balancing/firm_balancing.py}{39}{60}{Règlement des écarts, journée d'un producteur et prix capté.}\label{lst:m3:power-markets-trading:settle}
\item \textbf{La batterie.} \Cref{met:m3:power-markets-trading:battery}.
\omcode{code/markets-3/06-power-markets-trading/python/m3_trading.py}{41}{60}{Meilleur programme J-1 d'une batterie par programmation dynamique.}\label{lst:m3:power-markets-trading:dp}
\item \textbf{Exécuter} \texttt{m3\_trading.wind\_day()}, \texttt{capture(2025)},
\texttt{battery\_year(2025)} et \texttt{fig\_trading.py}.
\end{tutsteps}
\textbf{Pour aller plus loin.} Remplacer la prévision parfaite par les prix de la veille comme prévision
et mesurer le revenu perdu~; laisser la batterie proposer des produits de 15~minutes aux enchères à partir
d'octobre 2025.
\end{tutorial}

\section{Construction~: le gardien du groupe d'équilibrage}\label{bld:m3:power-markets-trading:balancing}

\begin{build}
\bfield{Rôle} La mini-entreprise exploite un groupe d'équilibrage pour ses actifs éoliens, solaires
et de batterie~: elle doit savoir, par MTU, ce qu'elle a programmé, ce qu'elle a mesuré, ce qu'elle doit
à l'opérateur du réseau et ce que sa production valait sur le marché au comptant.
\bfield{Interface} \texttt{BalancingGroup.trade(mtu, qty, price)},
\texttt{schedule\_production}, \texttt{schedule(mtu)}~; \texttt{imbalance}~;
\texttt{settle(volume, price\_long, price\_short)}~; \texttt{day\_settlement(sales, metered,
imbalance\_price)}~; \texttt{capture\_price(prices, output)}.
\bfield{Règles} Les quantités positives sont de l'énergie entrant dans le groupe~; un écart long est vendu à
l'opérateur, un écart court lui est acheté, à un prix unique sauf si deux prix sont donnés~; le \omterm{def:m3:power-markets-trading:capture}{prix capté} est
non défini sans production.
\bfield{Tests d'acceptation} \texttt{code/firm/balancing/tests/}~: un producteur équilibré~; les signes des
écarts et leurs flux financiers~; le règlement d'une journée~; un \omterm{def:m3:power-markets-trading:capture}{prix capté}.
\bfield{Pour aller plus loin} Des MTU d'un quart d'heure~; une tarification duale des écarts~; les nominations
du \cref{ch:m3:getting-access-commodities-and-power} construites à partir de ce gardien.
\end{build}

\begin{omsources}
\item ENTSO-E, Single Intraday Coupling~; Austrian Power Grid,
``SIDC: successful go-live of the 30-minute gate closure time in cross-zonal intraday trading''.
\item TSO allemands, \emph{Description of the balancing process and the balancing markets in
Germany}~; netztransparenz.de, prix de règlement des écarts uniforme (reBAP).
\item PJM, Financial Transmission Rights~; PJM Manual 06 et la FTR forfeiture rule.
\item Bundesnetzagentur | SMARD.de, prix horaires et production 2024--2025.
\item Pexapark, ``Unpacking H1 Deal Flow'' (2025)~; M.~Stover (TSSA), ``Understanding Battery Energy Storage Systems --- Current and Future'',
présentation devant le ERCOT Technology and Security Committee, 9 février 2026.
\end{omsources}

\section{Exercices}

\begin{exercise}[$\star$]\label{exo:m3:power-markets-trading:1}
Un parc solaire produit 10, 40, 40 et 10~MWh pendant quatre heures aux prix de 80, 20, 10 et
\qty{90}{EUR/MWh}. Calculer son \omterm{def:m3:power-markets-trading:capture}{prix capté}, le prix en base et son taux de capture.
\end{exercise}

\begin{exercise}[$\star$]\label{exo:m3:power-markets-trading:2}
Un \omterm{def:m3:power-markets-trading:brp}{responsable d'équilibre} (BRP) est en déficit de 20~MWh sur un quart d'heure dont le \omterm{def:m3:power-markets-trading:imbalance}{prix de règlement des écarts} est de \qty{300}{EUR/MWh}. Combien
paie-t-il~? Que recevrait-il s'il avait été en excédent de 20~MWh~?
\end{exercise}

\begin{exercise}[$\star$]\label{exo:m3:power-markets-trading:3}
Pourquoi existe-t-il un \omterm{def:m3:power-markets-trading:intraday}{marché infrajournalier}, si le \omterm{def:m3:power-markets-design:da}{marché J-1} fixe déjà un prix pour chaque
heure~?
\end{exercise}

\begin{exercise}[$\star\star$]\label{exo:m3:power-markets-trading:4}
Un \omterm{def:m3:power-markets-trading:ftr}{droit financier de transport} (FTR) de 50~MW du nœud $A$ au nœud $B$ est détenu pendant un jour. Les prix de congestion J-1 sont
\qty{-5}{} en $A$ et \qty{+15}{EUR/MWh} en $B$ pendant 10 heures, et égaux sinon. Combien
rapporte-t-il~?
\end{exercise}

\begin{exercise}[$\star\star$]\label{exo:m3:power-markets-trading:5}
Pour la batterie du texte, quel est le revenu d'une journée où les deux heures les moins chères coûtent
\qty{20}{} et les deux plus chères rapportent \qty{120}{EUR/MWh}, la plus chère venant après la moins chère~?
\end{exercise}

\begin{exercise}[$\star\star$]\label{exo:m3:power-markets-trading:6}
Un acheteur signe un \omterm{def:m3:power-markets-trading:ppa}{contrat d'achat d'électricité} (PPA) au fil de la production solaire à \qty{60}{EUR/MWh} pour 100~GWh par an. En prenant le
\omterm{def:m3:power-markets-trading:capture}{prix capté} 2025, combien ce PPA lui coûte-t-il au-delà de la valeur au comptant de la production~?
\end{exercise}

\begin{exercise}[$\star\star\star$]\label{exo:m3:power-markets-trading:7}
\emph{Programmation.} Avec \texttt{battery\_year}, calculer le revenu annuel par MW pour 2024 et 2025 avec
un cycle par jour, et pour 2025 avec deux.
\end{exercise}

\begin{exercise}[$\star\star\star$]\label{exo:m3:power-markets-trading:8}
\emph{Repérer la faille.} «~Notre backtest montre que la batterie gagne 75\,000~EUR par MW par an sur
les prix J-1~; c'est ce que nous gagnerons.~»
\end{exercise}

\section{Problème~: l'année d'une batterie}

\begin{problem}[{Problème du week-end --- 100 MW / 200 MWh sur le marché J-1 allemand}]\label{pb:m3:power-markets-trading:1}
Un investisseur envisage une batterie de 100~MW et 200~MWh, avec un \omterm{def:m3:power-markets-trading:rte}{rendement aller-retour} de 88~\%, autorisée
à effectuer un cycle complet par jour, n'échangeant que l'enchère J-1 allemande, avec une prévision parfaite.

\medskip\noindent\textbf{Partie I --- Une journée.}
\begin{enumerate}
\item Le 18 juin 2025, quand la batterie charge-t-elle et décharge-t-elle, et pourquoi à ce moment-là~?
\item Quelle énergie prélève-t-elle sur le réseau, et quelle énergie restitue-t-elle~?
\item Que gagne-t-elle ce jour-là~?
\item Pourquoi charge-t-elle à des prix négatifs plutôt qu'à zéro~?
\item Qu'apporterait un second cycle~?
\end{enumerate}

\medskip\noindent\textbf{Partie II --- L'année.}
\begin{enumerate}[resume]
\item Donner le revenu par MW en 2024 et en 2025.
\item Pourquoi 2025 a-t-elle été meilleure~?
\item Qu'apporte un second cycle quotidien en 2025~?
\item Pourquoi le revenu se concentre-t-il sur certains mois~?
\item Quel est, à peu près, le coût du rendement sur une journée comme le 18 juin~?
\end{enumerate}

\medskip\noindent\textbf{Partie III --- Réalisme.}
\begin{enumerate}[resume]
\item Pourquoi la prévision parfaite surestime-t-elle le revenu~?
\item Quels revenus le calcul laisse-t-il de côté~?
\item Que deviennent ces spreads lorsque de nombreuses batteries sont construites~?
\item Comment les cycles limitent-ils la durée de vie de la batterie, et pourquoi imposer un cap~?
\item Que changerait un marché au quart d'heure~?
\end{enumerate}

\medskip\noindent\textbf{Partie IV --- Jugement.}
\begin{enumerate}[resume]
\item Pourquoi les taux de capture du solaire et les revenus des batteries évoluent-ils ensemble~?
\item Qui d'autre se dispute les heures creuses de midi~?
\item Le spread J-1 est-il un revenu durable~?
\item Énoncer le \emph{résultat nommé}~: le revenu d'arbitrage J-1 d'une batterie par MW-année.
\item En une phrase~: que vend une batterie~?
\end{enumerate}
\end{problem}

\section{Questions d'entretien}

\begin{interviewq}[$\star$ \iqroles{trader}]\label{iq:m3:power-markets-trading:1}
Qu'est-ce qu'un \omterm{def:m3:power-markets-trading:imbalance}{prix de règlement des écarts}, et à quel moment être en position vendeuse revient-il bon marché~?
\end{interviewq}

\begin{interviewq}[$\star$ \iqroles{trader, chercheur}]\label{iq:m3:power-markets-trading:2}
Pourquoi le \omterm{def:m3:power-markets-trading:capture}{prix capté} du solaire est-il inférieur au prix de base, et pourquoi continue-t-il de baisser~?
\end{interviewq}

\begin{interviewq}[$\star\star$ \iqroles{chercheur}]\label{iq:m3:power-markets-trading:3}
Comment valoriseriez-vous un PPA éolien de dix ans payé selon la production~?
\end{interviewq}

\begin{interviewq}[$\star\star$ \iqroles{trader}]\label{iq:m3:power-markets-trading:4}
Expliquez ce qu'est une \omterm{def:m3:power-markets-trading:virtual}{offre virtuelle}, et comment elle pourrait être utilisée pour manipuler le marché.
\end{interviewq}

\begin{interviewq}[$\star\star$ \iqroles{risque}]\label{iq:m3:power-markets-trading:5}
L'erreur de prévision d'un portefeuille éolien atteint 10~\% de la production. Qu'est-ce qui détermine son
coût d'écart attendu~?
\end{interviewq}

\begin{interviewq}[$\star\star\star$ \iqroles{développeur, chercheur}]\label{iq:m3:power-markets-trading:6}
Concevez l'optimiseur qui place simultanément une offre de batterie sur l'enchère J-1, sur le \omterm{def:m3:power-markets-trading:intraday}{marché
infrajournalier} et sur un marché de réserve.
\end{interviewq}
